\documentclass[tikz,border=0pt]{standalone}
% Shared visual language for every block diagram in figures/diagrams.
% Build with XeLaTeX (see figures/diagrams/build.sh).
\usepackage{fontspec}
\setmainfont{Inter}[Numbers=Lining]
\setmonofont{DejaVu Sans Mono}[Scale=0.86]
\usepackage{xcolor}
\usetikzlibrary{arrows.meta,positioning,calc,fit,backgrounds,shapes.geometric,
  shapes.misc,matrix,decorations.pathreplacing,patterns,quotes}

% Palette (matches python/figstyle.py) --------------------------------------
\definecolor{surface}{HTML}{FCFCFB}
\definecolor{ink}{HTML}{0B0B0B}
\definecolor{ink2}{HTML}{52514E}
\definecolor{muted}{HTML}{898781}
\definecolor{hair}{HTML}{E1E0D9}
\definecolor{memline}{HTML}{2A78D6}  \definecolor{memfill}{HTML}{E7F0FC}
\definecolor{aluline}{HTML}{14875E}  \definecolor{alufill}{HTML}{E3F5EE}
\definecolor{ctlline}{HTML}{B07800}  \definecolor{ctlfill}{HTML}{FDF2D9}
\definecolor{regline}{HTML}{6250D6}  \definecolor{regfill}{HTML}{ECEAFB}
\definecolor{wire}{HTML}{3D3C39}
\definecolor{accent}{HTML}{EB6834}   \definecolor{accentfill}{HTML}{FDE8DF}
\definecolor{pos}{HTML}{E34948}      \definecolor{posfill}{HTML}{FBE1E0}
\definecolor{neg}{HTML}{2A78D6}      \definecolor{negfill}{HTML}{DCE9FA}
\definecolor{zerofill}{HTML}{F0EFEC}

\tikzset{
  every picture/.style={line cap=round, line join=round, font=\small, text=ink},
  % canvas
  canvas/.style={show background rectangle,
    background rectangle/.style={fill=surface, draw=hair, rounded corners=8pt, line width=0.6pt},
    inner frame xsep=16pt, inner frame ysep=14pt},
  ttl/.style={font=\bfseries\large, anchor=north west, inner sep=0pt, text=ink},
  sub/.style={font=\small, anchor=north west, inner sep=0pt, text=ink2},
  % blocks
  blk/.style={rounded corners=3pt, line width=0.8pt, align=center, minimum height=9mm,
    inner xsep=5pt, inner ysep=3pt},
  mem/.style={blk, draw=memline, fill=memfill},
  alu/.style={blk, draw=aluline, fill=alufill},
  ctl/.style={blk, draw=ctlline, fill=ctlfill},
  reg/.style={blk, draw=regline, fill=regfill, rounded corners=1pt},
  io/.style={blk, draw=muted, fill=white},
  ghost/.style={blk, draw=muted, dashed, fill=surface},
  % arithmetic primitives
  add/.style={circle, draw=aluline, fill=white, line width=0.8pt, minimum size=6.2mm,
    inner sep=0pt, font=\normalsize},
  shf/.style={rectangle, rounded corners=1.5pt, draw=aluline, fill=white, line width=0.7pt,
    font=\ttfamily\footnotesize, inner xsep=2.5pt, inner ysep=2pt, minimum height=5mm},
  gain/.style={isosceles triangle, isosceles triangle apex angle=60, draw=aluline, fill=white,
    line width=0.7pt, inner sep=1.2pt, font=\footnotesize},
  dn/.style={circle, draw=wire, fill=white, line width=0.7pt, minimum size=8mm, inner sep=0pt,
    font=\footnotesize},
  mux/.style={trapezium, trapezium angle=72, shape border rotate=270, draw=ctlline, fill=ctlfill,
    line width=0.8pt, minimum height=13mm, minimum width=5mm, inner sep=1pt, font=\scriptsize},
  net/.style={draw=memline, fill=memfill, rounded corners=1.5pt, font=\footnotesize, inner xsep=3pt,
    inner ysep=1.6pt, minimum width=5.5mm, minimum height=4.8mm},
  % wires
  sig/.style={-{Latex[length=2.2mm,width=1.6mm]}, draw=wire, line width=0.75pt},
  bus/.style={-{Latex[length=2.6mm,width=2mm]}, draw=wire, line width=1.5pt},
  ctrl/.style={-{Latex[length=1.8mm,width=1.4mm]}, draw=ctlline, line width=0.65pt, dash pattern=on 3pt off 2pt},
  nohead/.style={-, draw=wire, line width=0.75pt},
  dot/.style={circle, fill=wire, inner sep=0pt, minimum size=3.2pt},
  % labels
  lbl/.style={font=\footnotesize, text=ink2, fill=surface, inner sep=1.2pt},
  lblc/.style={font=\footnotesize, text=ink2, inner sep=1.2pt},
  note/.style={font=\footnotesize, text=ink2, align=left, inner sep=0pt},
  head/.style={font=\bfseries\small, text=ink, inner sep=0pt},
  code/.style={font=\ttfamily\footnotesize, text=ink},
  group/.style={draw=hair, fill=white, rounded corners=6pt, line width=0.8pt, inner sep=7pt},
  grouplabel/.style={font=\bfseries\footnotesize, text=muted, anchor=south west, inner sep=2pt},
}

% Layers: groups sit above the canvas background and below the main layer.
\pgfdeclarelayer{groups}
\pgfsetlayers{background,groups,main}

% A register: box with a clock wedge at the bottom edge.
\newcommand{\clk}[1]{\draw[regline, line width=0.6pt] ($(#1.south)+(-1.4mm,0)$) -- ($(#1.south)+(0,1.6mm)$) -- ($(#1.south)+(1.4mm,0)$);}
% Bus width tick: \bw{(point)}{width}
\newcommand{\bw}[2]{\draw[wire, line width=0.6pt] ($#1+(-0.9mm,-1.2mm)$) -- ($#1+(0.9mm,1.2mm)$); \node[font=\scriptsize, text=ink2, above=0.4mm, inner sep=0.5pt] at #1 {#2};}
\newcommand{\var}[1]{\textit{#1}}
\newcommand{\usub}[1]{\ifcase#1 ₀\or ₁\or ₂\or ₃\or ₄\or ₅\fi}
\AtBeginDocument{\hyphenpenalty=10000\exhyphenpenalty=10000}

\begin{document}
\begin{tikzpicture}[canvas]
  \node[ttl] at (0,0) {Halo-tiled direct 2-D engine (ARCH 4)};
  \node[sub] at (0,-0.62) {Copy a (2\var{T}+3)² raw patch with a global periodic halo, then emit \var{T}² quartets from the local window. Tiles live in coefficient space.};
  % datapath
  \node[mem, minimum width=24mm, minimum height=16mm] (fr) at (1.4,-2.9) {\texttt{frame}\\\footnotesize \var{N} × 8 bit};
  \node[ctl, minimum width=27mm, minimum height=14mm] (ld) at (5.0,-2.9) {halo loader\\\footnotesize 1 pixel / clock};
  \node[mem, minimum width=27mm, minimum height=16mm] (pt) at (8.8,-2.9) {\texttt{patch}\\\footnotesize \var{P}² × 8 bit};
  \node[alu, minimum width=30mm, minimum height=14mm] (dm) at (12.9,-2.9) {four masks\\\footnotesize LL, LH, HL, HH};
  \node[mem, minimum width=27mm, minimum height=16mm] (co) at (17.0,-2.9) {\texttt{coeff}\\\footnotesize quartet writes};
  \draw[bus] (fr) -- (ld);
  \draw[bus] (ld) -- node[lbl, above] {write} (pt);
  \draw[bus] (pt) -- node[lbl, above] {25 reads} (dm);
  \draw[bus] (dm) -- node[lbl, above] {4 writes} (co);
  \node[ctl, align=left, anchor=north west, minimum width=64mm] (c1) at (3.6,-4.35) {\textbf{TILE\_FILL}\quad \var{py}, \var{px} = 0…\var{P}−1 · \var{P}² clocks};
  \node[ctl, align=left, anchor=north west, minimum width=60mm] (c2) at (11.2,-4.35) {\textbf{TILE\_EMIT}\quad \var{py}, \var{px} = 0…\var{T}−1 · \var{T}² clocks};
  \draw[ctrl] (c1.east) -- (c2.west);
  \draw[ctrl] (c2.south) -- ++(0,-0.35) -| node[lbl, pos=0.25, below] {next tile: \var{tx} += \var{T}, then \var{ty} += \var{T} · \var{N}/(4\var{T}²) tiles} (c1.south);
  % geometry: T = 3, P = 9
  \def\g{0.42} \def\gx{0.6} \def\gy{-6.4}
  \node[head, anchor=south west] at (\gx,\gy+0.1) {Patch geometry, \var{T} = 3, \var{P} = 9};
  \foreach \r in {0,...,8} \foreach \q in {0,...,8} {
    \pgfmathtruncatemacro{\inner}{(\r>1)&&(\r<7)&&(\q>1)&&(\q<7)}
    \ifnum\inner=1 \def\fc{alufill}\else\def\fc{ctlfill}\fi
    \draw[draw=surface, line width=0.9pt, fill=\fc] (\gx+\q*\g,\gy-\r*\g) rectangle ++(\g,-\g);
  }
  \foreach \r in {0,1,2} \foreach \q in {0,1,2}
    \fill[aluline] (\gx+2*\q*\g+2.5*\g,\gy-2*\r*\g-2.5*\g) circle (2.4pt);
  \draw[accent, line width=1.2pt] (\gx+4*\g,\gy-4*\g) rectangle ++(5*\g,-5*\g);
  \node[note, anchor=north west, text width=6.3cm] at (\gx+9.4*\g,\gy-0.05) {%
    Gold: two-pixel halo on every side. Green: centre region.\\[2pt]
    Dots: the 3 × 3 local centres (2\var{py}+2, 2\var{px}+2), each producing one quartet.\\[2pt]
    Orange: the 5 × 5 window of the last centre; ±2 fits exactly in [0, 2\var{T}+2].\\[2pt]
    Neighbouring patches overlap by 3 rows or columns; the halo uses global wrapped coordinates, so tiles have no seams.};
  % numbers
  \node[note, anchor=north west, text width=8.2cm] at (11.6,\gy-0.05) {%
    \textbf{Cost per frame}\\[2pt]
    compute clocks = (\var{N}/4\var{T}²)((2\var{T}+3)² + \var{T}²)\\
    frame reads per pixel = (2\var{T}+3)²/(4\var{T}²)\\
    patch = 8\var{P}² bits\\[4pt]
    \begin{tabular}{@{}rrrr@{}}
      \var{T} & clocks/\var{N} & reads/pixel & patch\\
      1 & 6.50 & 6.25 & 200 b\\
      4 & 2.14 & 1.89 & 968 b\\
      8 & 1.66 & 1.41 & 2,888 b\\
      32 & 1.35 & 1.10 & 35,912 b\\
    \end{tabular}\\[3pt]
    Large tiles approach 1.25\var{N} clocks and one read per pixel.};
\end{tikzpicture}
\end{document}
